\documentclass[10pt,a4paper]{amsart}
\usepackage[margin=19mm]{geometry}
\usepackage{amsmath,amssymb,mathtools}
\usepackage[scr=rsfs,cal=euler]{mathalfa}
\usepackage{xfrac}
\usepackage[unicode,hidelinks,bookmarks=false]{hyperref}
\newcommand{\ie}{i.e.\ }
\newcommand{\cf}{cf.\ }
\newcommand{\resp}{respectively\ }
\newcommand{\Subsection}{\subsection}
\providecommand{\nobreakdash}{\nobreak\hbox{-}}
\setlength{\emergencystretch}{2em}
\multlinegap0pt
\usepackage{lmodern}
\usepackage{mathtools}
\usepackage{booktabs}
\usepackage{mathrsfs}
\newtheorem{theorem}{Theorem}[section]
\newtheorem{proposition}[theorem]{Proposition}
\newtheorem{lemma}[theorem]{Lemma}
\newtheorem{corollary}[theorem]{Corollary}
\newtheorem{remark}[theorem]{Remark}
\newtheorem{conjecture}[theorem]{Conjecture}
\newtheorem{problem}[theorem]{Problem}
\newcommand{\mdr}{\operatorname{mdr}}
\newcommand{\AR}{\operatorname{AR}}
\newcommand{\ii}{\mathrm i}
\newcommand{\Sing}{\operatorname{Sing}}
\begin{document}
\title[On the Numerical Terao Conjecture]{On the Numerical Terao Conjecture}
\author{Piotr Pokora}
\date{}
\maketitle
\noindent\emph{Source and licence.} On the Numerical Terao Conjecture. Version arXiv:2608.05662v3 (CC BY 4.0). This material is distributed under the Creative Commons Attribution 4.0 International licence, http://creativecommons.org/licenses/by/4.0/, as recorded in the arXiv OAI licence metadata for this exact version and shown on the abstract page https://arxiv.org/abs/2608.05662v3. Full rights notice: licenses/arxiv-2608.05662-CC-BY-4.0.txt.\par\smallskip
\noindent\textit{Editorial scope.} Counted unit: the original \texttt{theorem} environment \texttt{thm:counterexample} at source lines 267--277 together with its complete proof at lines 278--373; the delivered document keeps source lines 74--380 so that every referenced label and every citation resolves inside it. Native editable source is the paper's own concatenated TeX; the AMS wrapper is editorial formatting only. No publisher class, upstream script or unrelated artwork is redistributed.
\begin{equation}\label{eq:DS}
 \mdr(f)\ge \alpha_{\mathcal{C}} \cdot d-2.
\end{equation}
We also use the following du Plessis-Wall upper bound \cite{DPP1} and its freeness characterization in the form recalled in~\cite{DimcaMaxTau,DimcaSticlaru}.
For a reduced plane curve $\mathcal{C} : f=0$ we put $r=\mdr(f)$.  If $r<d/2$, then
\begin{equation}\label{eq:DPW-small}
 \tau(\mathcal{C})\le (d-1)^2-r(d-r-1),
\end{equation}
and equality holds precisely when $\mathcal{C}$ is free.  If $r\ge d/2$, then the
refined bound gives us
\begin{equation}\label{eq:DPW-large}
 \tau(\mathcal{C})\le (d-1)^2-r(d-r-1) -\binom{2r-d+2}{2}.
\end{equation}

\section{A structure theorem on conics and lines with ADE singularities}

\begin{lemma}\label{lem:types}
Let $\mathcal{CL}$ be a reduced plane arrangement of lines and smooth conics having only ADE singularities.  Then every singular point is of one of the following types:
\[
 A_1,A_3,A_5,A_7,\qquad D_4,D_6,D_8,D_{10}.
\]
Consequently,
\[
 \alpha_{\mathcal{CL}}\ge \alpha(D_{10})=\frac59.
\]
\end{lemma}

\begin{proof}
Every local branch of $\mathcal{CL}$ is smooth.  Hence an $A$-type singularity
which can occur in the arrangement has two smooth branches and is of type
$A_{2q-1}$.  The number $q$ is their local intersection multiplicity.  Since
the components have degrees at most two, B\'ezout's theorem gives $q\le4$.
Thus only $A_1,A_3,A_5,A_7$ can occur.

An ADE singularity with three smooth branches is necessarily of type
$D_{2q+2}$: two branches have contact order $q$, and a third branch is
transverse to them.  Again $q\le4$, so the possible types are
$D_4,D_6,D_8,D_{10}$.  The remaining simple singularities have either a
singular local branch or a different branch structure and therefore cannot occur in a union of smooth components. Let us recall from \cite{DimcaSernesi} that for our quasi-homogeneous singularities one has
\[
 \alpha(A_k)=\frac12+\frac1{k+1},
 \qquad
 \alpha(D_k)=\frac{k}{2(k-1)}.
\]
Among the listed types the smallest value is
\[
 \alpha(D_{10})=\frac{10}{18}=\frac59.
\]
\end{proof}

\begin{theorem}\label{thm:even}
Let $\mathcal{CL} = \{f=0\} \subset \mathbb{P}^{2}_{\mathbb{C}}$ be a conic-line arrangement of even degree
$d=2m\ge4$ with only ADE singularities.  Suppose that $\mathcal{CL}$ is free. Then every conic-line arrangement with the same weak combinatorics is free. Consequently, there is no counterexample to the Numerical Terao Conjecture in even degree in the ADE conic-line class.
\end{theorem}

\begin{proof}
Set $r=\mdr(f)$.  Since $\mathcal{CL}$ is free of degree $2m$, its exponents
are $(r,2m-1-r)$, and therefore
\begin{equation}\label{eq:r-upper}
 r\le m-1.
\end{equation}
By Lemma~\ref{lem:types} and \eqref{eq:DS},
\begin{equation}\label{eq:r-lower}
 r\ge \left\lceil\frac{10m}{9}-2\right\rceil.
\end{equation}
For $2\le m\le9$ one has
\[
 \left\lceil\frac{10m}{9}-2\right\rceil=m-1.
\]
Combining this with \eqref{eq:r-upper} shows that
\begin{equation}\label{eq:middle-exp}
 r=m-1,
 \qquad
 \exp(\mathcal C)=(m-1,m).
\end{equation}
If $m\ge10$, then the right-hand side of \eqref{eq:r-lower} is at least
$m$, contradicting \eqref{eq:r-upper}.  Hence no free ADE conic-line
arrangement of even degree $d\ge20$ exists.

It remains to prove combinatorial rigidity in the possible degrees
$4\le d\le18$.  Let $\mathcal C':g=0$ have the same weak combinatorics as
$\mathcal C$, and put $r'=\mdr(g)$.  The two curves have the same total
Tjurina number.  By \eqref{eq:middle-exp},
\begin{equation}\label{eq:tau-even-free}
 \tau(\mathcal C')=\tau(\mathcal C)
 =(2m-1)^2-m(m-1).
\end{equation}
The universal Arnold-exponent estimate again gives $r'\ge m-1$. If $r'=m-1$, then \eqref{eq:tau-even-free} is equality in
\eqref{eq:DPW-small}, and therefore $\mathcal C'$ is free.

Assume now that $r'\ge m$, and write $r'=m+s$ with $s\ge0$.  Applying the
refined bound \eqref{eq:DPW-large}, we obtain
\begin{align*}
 \tau(\mathcal C')
 &\le (2m-1)^2-(m+s)(m-1-s)-\binom{2s+2}{2}\\
 &= (2m-1)^2-m(m-1)-(s+1)^2\\
 &=\tau(\mathcal C)-(s+1)^2,
\end{align*}
which contradicts \eqref{eq:tau-even-free}.  Thus the second case is
impossible, and $\mathcal C'$ is necessarily free.
\end{proof}
The above result gives us a surprising boundedness result on maximizing conic-line arrangements. By \cite[Theorem 2.9]{max}, an ADE reduced plane curve of even degree $d=2m$ is maximizing if and only if it is free with exponents
\[
(m-1,m).
\]
\begin{corollary}
Let $\mathcal{CL}\subset \mathbb{P}^2_{\mathbb{C}}$ be a conic-line arrangement of even degree
\[
d=2m\geq 4
\]
having only ADE singularities. If $\mathcal{CL}$ is maximizing, then
\[
d\leq 18.
\]
\end{corollary}

\begin{proof}
Every maximizing conic-line arrangement is a free
ADE conic-line arrangement. By Theorem \ref{thm:even}, free conic-line arrangements with only ADE singularities are bounded in degree, and necessarily
\[
d\leq 18.
\]
\end{proof}
Finishing this section, we would like to focus on the odd degree case.
\begin{proposition}
\label{odd}
Let $\mathcal{CL} = \{f=0\} \subset \mathbb{P}^{2}_{\mathbb{C}}$ be a conic-line arrangement of odd degree
$$d = 2m+1 \geq 5$$
having only ${\rm ADE}$ singularities. If $\mathcal{CL}$ is free, then
$$d \leq 27.$$
\end{proposition}
\begin{proof}
Let $r=\operatorname{mdr}(f)$. Since $\mathcal{CL}$ is free of degree
$d=2m+1$, its exponents are $(r,2m-r)$, and hence $r\leq m$.
By Lemma~\ref{lem:types} and the Dimca--Sernesi bound, we have
\[
m\geq r\geq \alpha_{\mathcal{CL}}\cdot d-2
\geq \frac{5}{9}(2m+1)-2
=\frac{10m}{9}-\frac{13}{9}.
\]
Consequently $m\leq 13$, and therefore $d=2m+1\leq 27$.
\end{proof}
Let us recall that by \cite[Proposition 5.1(a)]{max}, an ADE reduced plane curve of odd degree $d=2m+1 \geq 5$ is maximizing if and only if it is free with exponents
\[
(m-1,m+1).
\]
\begin{corollary}
Let $\mathcal{CL}\subset \mathbb{P}^2_{\mathbb{C}}$ be a conic-line arrangement of odd degree
\[
d=2m+1\geq 5
\]
having only ADE singularities. If $\mathcal{CL}$ is maximizing, then
\[
d \in \{5,7,9\}.
\]
\end{corollary}
\begin{proof}
Let $r = {\rm mdr}(f)$. Since $\alpha_{\mathcal{CL}} \geq \frac{5}{9}$, we have
$$ m-1 = r \geq \alpha_{\mathcal{CL}}\cdot d -2  \geq  \frac{10m}{9} - \frac{13}{9},$$
hence $m \leq 4$.
\end{proof}
Finally, combining Theorem \ref{thm:even} and Proposition \ref{odd}, we arrive at the following observation.
\begin{corollary}
There are only finitely many weak-combinatorial types of free conic-line arrangements with ${\rm ADE}$ singularities.
\end{corollary}
\begin{proof}
We have already observed that the degree of a free conic-line
arrangement with only ADE singularities is bounded by $27$. Moreover,
by Lemma~\ref{lem:types}, the only possible singularity types are
\[
A_1,\ A_3,\ A_5,\ A_7,\ D_4,\ D_6,\ D_8,\ D_{10}.
\]
For a fixed total degree $d\leq 27$, there are only finitely many
possibilities for the numbers of line and conic components. Furthermore,
the number of singular points of each of the above types is bounded,
for instance by the total intersection number of the irreducible
components. Consequently, only finitely many weak-combinatorial types
can occur.
\end{proof}
\section{On a special pair of conic-line arrangements}

Put $S=\mathbb C[x,y,z]$ and let $\ii^2=-1$.  Define
\begin{align*}
 F={}&y(y-x)(y+x)(y-x-2z)(y-x+2z)\\
 &\qquad\times (y+x-2z)(y+x+2z)
       (x^2-y^2+4\ii yz-4z^2),
\end{align*}
and
\begin{align*}
 G={}&xyz(x-y)(2x-z)(y-z)(2x+y-z)\\
 &\qquad\times \bigl(z^2-yz-2xz+(1-\ii)xy\bigr).
\end{align*}
We set $\mathcal C_F : F=0$ and $\mathcal C_G : G=0$. Obviously the two quadratic factors define smooth conics. 

\begin{theorem}\label{thm:counterexample}
The curves $\mathcal C_F$ and $\mathcal C_G$ have the same weak
combinatorics, namely
\[
 W(\mathcal C_F)=W(\mathcal C_G)
 =(7,1;\,8A_1+D_4+4X_9).
\]
Moreover, $\mathcal C_F$ is free with exponents $(4,4)$, whereas
$ \mdr(G)=3$ and $\mathcal{C}_G$ is nearly free with exponents $(3,6)$.
Thus the pair is a counterexample to the Numerical Terao Conjecture in the class of conic-line arrangements with ordinary quasi-homogeneous singularities. 
\end{theorem}

\begin{proof}
The seven-line part of $F$ has six ordinary double points and the following five ordinary triple points:
\[
 (0:0:1),\quad
 (1:0:-\tfrac12),\quad
 (1:0:\tfrac12),\quad
 (1:1:0),\quad
 (1:-1:0).
\]
Its conic passes transversely through the last four points and it avoids the first one and all six double points. Thus the first point remains a
$D_4$-point, whereas the other four triple points become ordinary quadruple
points.  The only unused conic-line intersection multiplicities occur on
$y-x=0$ and $y+x=0$; their residual points are
\[
 (1:1:\ii),\qquad (1:-1:-\ii),
\]
and both are transverse.

The seven-line part of $G$ likewise has six ordinary double points and five ordinary triple points. Four of the triple points are
\[
 (0:1:0),\quad (0:1:1),\quad (1:0:0),\quad (1:0:2),
\]
and the conic passes transversely through all the four. The remaining triple point is $(0:0:1)$, which the conic avoids. The only line with unused conic-intersection degree is $x-y=0$.  Restriction of the conic to this line gives
\[
 z^2-3yz+(1-\ii)y^2,
\]
whose discriminant is $5+4\ii\ne0$.  Hence it contributes two distinct
transverse nodes. It follows in both cases that
\[
 n_2=8,\qquad n_3=1,\qquad n_4=4,
\]
or, in singularity notation according to Arnold's catalogue \cite{arnold},
\[
 8A_1+D_4+4X_9.
\]
The total Tjurina number is therefore
\begin{equation}\label{eq:tau-counterexample}
 \tau=8\cdot1+1\cdot4+4\cdot9=48.
\end{equation}
All singularities of the two curves are quasi-homogeneous and the smallest Arnold exponent is now the value $1/2$ which corresponds to $X_9$. Hence
\begin{equation}\label{eq:lower-X9}
 \mdr(H)\ge 9\cdot\frac12-2=\frac52,
\text{ and it implies that }\mdr(H)\ge3
\end{equation}
for $H \in \{F,G\}$.

Now we observe that there is an explicit cubic Jacobian relation for $G$, and this computation can be performed using \texttt{SINGULAR}. Set
\begin{align*}
 a={}&-x(4x-5y)(2x+y-2z),\\
 b={}& y(5x-4y)(2x+y-2z),\\
 c={}& z(10x^2-12xy-xz+5y^2-yz).
\end{align*}
A direct check shows that
\[
 aG_x+bG_y+cG_z=0.
\]
Together with \eqref{eq:lower-X9}, this proves
\[
 \mdr(G)=3.
\]
For $d=9$ and $r=3$, the maximal Tjurina number is
\[
 \tau_{\max}(9,3)=8^2-3(8-3)=49,
\]
and such a curve attaining the maximal values of $\tau_{\max}(9,3)$ would be an $\mathscr{M}$-arrangement of conics and lines in the sense of \cite{MJ}. By \eqref{eq:tau-counterexample} we have
\[
 \tau(\mathcal C_G)=48=\tau_{\max}(9,3)-1.
\]
Using \cite{DimcaMaxTau} we can conclude that $\mathcal C_G$ is nearly free with exponents $(3,6)$, and in particular is not free.

Let us look at $\mathcal{C}_{F}$. Using \texttt{SINGULAR}, we can directly check that $\AR(F)_3=0$, and this means that
\[
 \mdr(F)\geq 4.
\]
We set 
\[
\begin{aligned}
a &= x^{4}-10x^{2}y^{2}+9y^{4}
     +28x^{2}z^{2}-36y^{2}z^{2},\\
b &= -8x^{3}y+8xy^{3}-8xyz^{2},\\
c &= 10x^{3}z-10xy^{2}z-8xz^{3}.
\end{aligned}
\]
These polynomials satisfy the Jacobian syzygy

\[
aF_x+bF_y+cF_z=0,
\]
hence ${\rm mdr}(F)=4$.
Finally,
\[
 \tau(\mathcal C_F)=48=8^2-4(8-4),
\]
so equality holds in \eqref{eq:DPW-small}.  It follows that $\mathcal C_F$
is free with exponents $(4,4)$.
\end{proof}

\begin{remark}
The existence of a free conic-line arrangement with weak combinatorics
\[
W(\mathcal{C})=(7,1;8A_1+D_4+4X_9)
\]
was already established in \cite[Theorem~1.3]{PokoraOneConic}. For the reader's convenience, we present here a substantially simpler defining equation.
\end{remark}

\begin{thebibliography}{99}
\bibitem[DPP1]{DPP1} A. A. du Plessis and C. T. C. Wall, Application of the theory of the discriminant to highly singular plane curves. \textit{Math. Proc. Camb. Philos. Soc.} \textbf{126(2)}: 259 -- 266 (1999).
\bibitem[DimcaMaxTau]{DimcaMaxTau} A.~Dimca, Freeness versus maximal global Tjurina number for plane curves. \textit{Math. Proc. Cambridge Philos. Soc.} \textbf{163}: 161--172 (2017).
\bibitem[DimcaSernesi]{DimcaSernesi} A.~Dimca and E.~Sernesi, Syzygies and logarithmic vector fields along plane curves. \textit{J. \`Ec. polytech. Math.} \textbf{1}: 247--267 (2014).
\bibitem[DimcaSticlaru]{DimcaSticlaru} A.~Dimca and G.~Sticlaru, Jacobian syzygies, fitting ideals, and plane curves with maximal global Tjurina numbers. \textit{Collect. Math.} \textbf{73(3)}: 391 -- 409 (2022).
\bibitem[MJ]{MJ} M. Janasz and P. Pokora, On $\mathscr{M}$-arrangements of conics and lines with ordinary singularities. \textbf{arXiv:2512.24707} (2026).
\bibitem[PokoraOneConic]{PokoraOneConic} P.~Pokora, Freeness of arrangements of lines and one conic with ordinary quasi-homogeneous singularities. \textit{Taiwanese J. Math.} \textbf{29(6)}: 1651 -- 1666 (2025).
\bibitem[arnold]{arnold} V. I. Arnold, Local normal forms of functions. \textit{Invent. Math.} \textbf{35}: 87 -- 109 (1976).
\bibitem[max]{max} A. Dimca and P. Pokora, Maximizing curves viewed as free curves. \textit{Int. Math. Res. Not.} \textbf{2023(22)}: 19156 -- 19183 (2023).
\end{thebibliography}
\end{document}
